\centerline{\bf \Huge Круги Эйлера II}



\problem{\textbf{1}} В классе 29 человек. 15 из них занимаются в музыкальном кружке, 21 — в математическом. Сколько человек посещают оба кружка, если известно, что только Вовочка(ученик класса) не ходит ни в один из двух кружков?


\problem{\textbf{2}} У Пети есть карточки, у которых каждая сторона — красная или синяя.
Карточек с синей стороной — 10, карточек с красной стороной — 12, карточек
с разными сторонами — 5. Сколько карточек, у которых стороны одного цвета? 

\problem{\textbf{3}} В киоске около школы продается мороженое двух видов:
«Спортивное» и «Мальвина». На перемене 24 ученика успели купить мороженое. При этом 15 из них купили «Спортивное», а 17 — мороженое «Мальвина». Сколько человек купили мороженое обоих сортов?

\begin{flushright} \textbf{2 балла}\\ 
\end{flushright}


\problem{\textbf{4}} В летнем лагере 70 ребят. Из них 27 занимаются в драмкружке,
32 поют в хоре, 22 увлекаются спортом. В драмкружке 10 ребят из хора, в хоре 6 спортсменов,
в драмкружке 8 спортсменов; 3 спортсмена посещают и драмкружок, и хор. Сколько ребят не
поют в хоре, не увлекаются спортом и не занимаются в драмкружке?

\problem{\textbf{5}} 
На олимпиаде было предложено две задачи. После проверки жюри составило четыре списка. В первый список внесли решивших первую задачу (включая и тех, кто решил обе задачи), во второй — решивших только одну задачу, а в третий — решивших по меньшей мере одну задачу. Четвёртый же список состоял из решивших обе задачи.

а) Какой из списков самый длинный?

б) Могут ли какие–то три списка совпадать по составу? Если да, то какие? А все четыре — могут?


\problem{\textbf{6}} На пиратском корабле трудятся 67 морских разбойников. У 47 из них есть
ухо, у 35 — глаз, а у 23 счастливчиков есть и то, и другое. Новая инструкция
Профсоюза Работников Абордажного Крюка предписывает корабельному врачу
учитывать также наличие носа. Оказалось, что 20 пиратов имеют нос, 12 — и
нос, и ухо, 11 — и нос, и глаз, а 5 — все три органа. Сколько пиратов не имеют
ничего из перечисленного?

\begin{flushright} \textbf{4 балла}\\ 
\end{flushright}

\hspace{3mm}

\problem{\textbf{7}} На прогулку пошли шестиклассники и пятиклассники. Все они были либо
босиком, либо в тапочках. Шестиклассников было 24, а босых учеников 16.
Обутых пятиклассников было столько же, сколько босых шестиклассников.
Сколько учеников ходили на прогулку? 

\problem{\textbf{8}} В кружок робототехники берут только тех, кто знает математику, физику или программирование. Известно, что 8 членов кружка
знают физику, 7 — математику, 11 — программирование. При этом известно, что не менее двоих
знают одновременно физику и математику, не менее троих — математику и программирование,
и не менее четырёх — физику и программирование. Какое наибольшее количество участников
кружка может быть при этих условиях?

\begin{flushright} \textbf{6 баллов}\\ 
\end{flushright}